\chapter{Referência Rápida}
\label{ap:referencia}

Este apêndice reúne, em formato de consulta rápida, tabelas já apresentadas ao longo da apostila — útil como cola de referência durante os primeiros projetos, sem precisar percorrer os capítulos inteiros novamente.

\section{Precedência de operadores}

Da maior para a menor precedência (operadores na mesma linha têm precedência igual, avaliados conforme sua associatividade):

\begin{table}[h]
\centering
\begin{tabular}{@{}ll@{}}
\toprule
\textbf{Operadores} & \textbf{Categoria} \\
\midrule
\code{()} \ \code{[]} \ \code{->} \ \code{.}        & chamada, indexação, acesso a membro \\
\code{!} \ \code{\textasciitilde} \ \code{++} \ \code{--} \ \code{(tipo)} \ \code{*} \ \code{\&} \ \code{sizeof} & unários, cast, ponteiro \\
\code{*} \ \code{/} \ \code{\%}                     & multiplicativos \\
\code{+} \ \code{-}                                 & aditivos \\
\code{<<} \ \code{>>}                               & deslocamento de bits \\
\code{<} \ \code{<=} \ \code{>} \ \code{>=}         & relacionais \\
\code{==} \ \code{!=}                               & igualdade \\
\code{\&}                                            & E bit a bit \\
\code{\textasciicircum}                             & OU exclusivo bit a bit \\
\code{|}                                             & OU bit a bit \\
\code{\&\&}                                          & E lógico \\
\code{||}                                            & OU lógico \\
\code{?:}                                            & ternário \\
\code{=} \ \code{+=} \ \code{-=} \ \code{*=} \ldots  & atribuição \\
\code{,}                                             & vírgula \\
\bottomrule
\end{tabular}
\caption{Precedência de operadores em C, do maior para o menor.}
\end{table}

\begin{dica}
Na dúvida sobre a ordem de avaliação de uma expressão complexa, use parênteses explícitos — isso nunca está errado, e deixa a intenção do código clara tanto para o compilador quanto para quem lê depois.
\end{dica}

\section{Especificadores de formato (\code{printf}/\code{scanf})}

\begin{table}[h]
\centering
\begin{tabular}{@{}ll@{}}
\toprule
\textbf{Especificador} & \textbf{Tipo} \\
\midrule
\code{\%d} / \code{\%i} & \code{int} \\
\code{\%u}              & \code{unsigned int} \\
\code{\%ld} / \code{\%lu} & \code{long} / \code{unsigned long} \\
\code{\%f}              & \code{float} / \code{double} \\
\code{\%.Nf}            & ponto flutuante com N casas decimais \\
\code{\%c}              & \code{char} \\
\code{\%s}              & string (\code{char *}) \\
\code{\%x} / \code{\%X} & hexadecimal (minúsculo / maiúsculo) \\
\code{\%o}              & octal \\
\code{\%p}              & endereço de ponteiro \\
\code{\%zu}             & \code{size\_t} (retorno de \code{sizeof}) \\
\bottomrule
\end{tabular}
\caption{Especificadores de formato mais usados.}
\end{table}

\section{Tipos de largura fixa (\code{stdint.h})}

\begin{table}[h]
\centering
\begin{tabular}{@{}lc@{}}
\toprule
\textbf{Tipo} & \textbf{Intervalo} \\
\midrule
\code{uint8\_t}  & 0 a 255 \\
\code{int8\_t}   & $-128$ a 127 \\
\code{uint16\_t} & 0 a 65\,535 \\
\code{int16\_t}  & $-32\,768$ a 32\,767 \\
\code{uint32\_t} & 0 a $2^{32}\!-\!1$ \\
\code{int32\_t}  & $-2^{31}$ a $2^{31}\!-\!1$ \\
\bottomrule
\end{tabular}
\caption{Tipos inteiros de largura fixa.}
\end{table}

\section{Funções essenciais do framework Arduino usadas nesta apostila}

\begin{table}[h]
\centering
\begin{tabular}{@{}ll@{}}
\toprule
\textbf{Função / macro} & \textbf{Finalidade} \\
\midrule
\code{pinMode(pino, modo)}     & define um pino como \code{OUTPUT}, \code{INPUT} ou \code{INPUT\_PULLUP} \\
\code{digitalWrite(pino, nivel)} & escreve nível lógico (\code{HIGH}/\code{LOW}) em um pino de saída \\
\code{digitalRead(pino)}       & lê o nível lógico de um pino de entrada \\
\code{analogRead(pino)}        & realiza uma leitura do ADC (0 a 4095 no \ESP{}) \\
\code{delay(ms)}                & pausa cooperativa, em milissegundos \\
\code{millis()}                 & milissegundos desde a inicialização (para temporização sem bloqueio) \\
\code{Serial.begin(baud)}       & inicia a comunicação serial na velocidade indicada \\
\code{Serial.print} / \code{Serial.println} & escreve na porta serial, com ou sem quebra de linha \\
\code{Serial.printf(fmt, \ldots)} & escrita formatada na porta serial (especificadores de \code{printf}) \\
\code{attachInterrupt(pino, isr, modo)} & associa uma função ISR a um pino (\code{RISING}/\code{FALLING}/\code{CHANGE}) \\
\code{digitalPinToInterrupt(pino)} & converte um número de pino no identificador de interrupção \\
\code{IRAM\_ATTR}               & garante que uma ISR fique carregada na RAM interna \\
\code{log\_i} / \code{log\_w} / \code{log\_e} & mensagens de log por nível de severidade \\
\code{randomSeed} / \code{random} & inicializa a semente e gera números pseudoaleatórios \\
\code{esp\_task\_wdt\_init/add/reset} & configura, registra e alimenta o Watchdog Timer \\
\code{esp\_sleep\_enable\_timer\_wakeup} & configura um temporizador como fonte de despertar \\
\code{esp\_light\_sleep\_start()}      & entra em \emph{light sleep} (RAM preservada, execução retomada) \\
\code{esp\_sleep\_enable\_gpio\_wakeup()} / \code{gpio\_wakeup\_enable} & configura um pino como fonte de despertar do \emph{light sleep} \\
\code{esp\_deep\_sleep\_start()}       & coloca o \ESP{} em modo \emph{deep sleep} \\
\code{RTC\_DATA\_ATTR}                 & preserva uma variável entre ciclos de \emph{deep sleep} \\
\code{WiFi.begin(ssid, senha)}         & inicia a conexão a uma rede Wi-Fi \\
\code{WiFi.status()} / \code{WiFi.localIP()} & consulta o estado da conexão / o IP obtido \\
\code{HTTPClient} (\code{begin}/\code{POST}/\code{end}) & realiza requisições HTTP \\
\code{ArduinoOTA.begin()} / \code{.handle()} & habilita e processa atualizações de firmware pela rede \\
\bottomrule
\end{tabular}
\caption{Funções do framework Arduino apresentadas na Parte II desta apostila.}
\end{table}
